%!TEX TS-program = xelatex
%!TEX encoding = UTF-8 Unicode

%-------------------------------------------------------------------------------
% VSSS 2026 Cover Letter
% Same font/style as CV, without personal header
% Compile with XeLaTeX
% Put this file in the same folder as cv_v4.tex
%-------------------------------------------------------------------------------

\documentclass[11pt, a4paper]{src/awesome-cv}

\usepackage{setspace}

%-------------------------------------------------------------------------------
% GLOBAL SETTINGS
%-------------------------------------------------------------------------------
\raggedbottom
\setlength{\parskip}{0pt}
\setlength{\parindent}{0pt}
\setstretch{1}

%-------------------------------------------------------------------------------
% COLOR SETTING - same as CV
%-------------------------------------------------------------------------------
\definecolor{longblue}{HTML}{101CA4}
\colorlet{awesome}{longblue}
\setbool{acvSectionColorHighlight}{false}

%-------------------------------------------------------------------------------
\begin{document}

%-------------------------------------------------------------------------------
% FOOTER
%-------------------------------------------------------------------------------
\makecvfooter
  {}
  {Long Nguyen Khac~~~·~~~Cover Letter}
  {\thepage}

%-------------------------------------------------------------------------------
% COVER LETTER CONTENT
%-------------------------------------------------------------------------------

\begin{flushright}
{\today}
\end{flushright}

\vspace{0.35cm}

Dear Vietnam Summer School of Science (VSSS) 2026,

\vspace{0.15cm}

I did not come to science by a straight path from the beginning. My academic journey has been more like an adaptive control process: learning from errors, adjusting by feedback, and finding a more stable balance. For this reason, the theme of VSSS this year, \textbf{“Rethinking Science”} strongly reflects the period I am going through now: looking back how I understand science, how I do research, and how I prepare for a long-term academic path.

\vspace{0.15cm}

My interest in science started early, through scientific research competitions during my school years. At that time, what attracted me most was the feeling of turning an idea into a real model. I used to think that doing science mainly meant developing or improving an existing research topic. However, after joining a research lab at university, I gradually realized that the research is not only building a product or applying a method. More importantly, it is about understanding what the real problem is, where the current limits are, how can bridge these limits, and how a topic contributes to a larger academic picture.

\vspace{0.15cm}

During my university years, I had the opportunity to work on different research projects in electronics, artificial intelligence, and robotics to control systems. These experiences helped me build a technical and research foundation, but they also made me ask myself an important question: am I really building a clear academic direction, or only moving through separate experiences? This question became stronger when I entered my final year and started to think more seriously about my path toward graduate studies.

\vspace{0.15cm}

To test my direction more seriously, I decided to take a one-year gap period and step out of my comfort zone through international research experiences. In Taiwan, when working on cutting force and tool wear prediction in CNC machining, I realized the gap between a model and a real system. Input data is not always clean, measurement signals are not always ideal, and a model only becomes meaningful when it can help improve the system efficiency. This experience changed the way I saw research. It should not be separated from practical problems, but should answer clearer questions: What problem does it solve? Who does it serve? And why does it matter?

\vspace{0.15cm}

Later, in Singapore, when working on wireless power transfer systems for autonomous underwater robots, I received another important “feedback signal”. I realized that a modern robotic system cannot be understood from only one field. Mechanical design, control, communication, and energy systems are closely connected. This experience taught me that engineering research requires not only technical depth, but also the ability to view a problem as a whole system.

\vspace{0.15cm}

These international experiences also helped me see how fast the global academic world is changing. New research directions appear continuously, and young people now have more opportunities to enter leading academic environments. However, I also understand that these opportunities require a strong foundation, a serious research mindset, and a clear academic plan. This has become an important motivation for me to complete my current study program in Vietnam, build a meaningful graduate thesis, and prepare better for my future graduate studies.

\vspace{0.15cm}

My story with VSSS began when I was introduced to the program by my high school research supervisor. I applied before, but at that time, I may not have fully understood what I needed from an academic community, VSSS. After meeting some VSSS alumni in Singapore and listening to their stories after the summer school, I understood more clearly the value of VSSS. To me, VSSS is not only a summer school with inspiring lectures, but also an academic network where young Vietnamese researchers can connect, rethink, and gain more motivation to go further.

\vspace{0.15cm}

At this point in my journey,  I believe VSSS 2026 comes at the right time for me. I hope to join the program to reflect more deeply on my academic direction, and understand more clearly the responsibility of a young Vietnamese researcher who wants to pursue science and contribute to the development of science and technology in Vietnam. I also hope to listen to the stories and sharing of leading professors, scientists, and young people from different fields.
\vspace{0.15cm}

The image of harmonograph in this year’s theme gave me many thoughts. Each pendulum has its own movement, but the interaction between them can create a larger and more beautiful pattern. I really hope to become one small pendulum in that network, learning from those who came before me, connecting with people who share the same passion for science, and contributing my own small movement to the shared values that VSSS is spreading.

\vspace{0.15cm}

Sincerely,\\[0.15cm]
\textbf{Long Nguyen Khac}

%-------------------------------------------------------------------------------
\end{document}